\section*{第 \ref{ch:g12:plant-rooted-life} 章——植物生命：扎根而欣欣向荣}

\begin{solution}{exo:g12:plant-rooted-life:1}
根系，其\omterm{prop:g12:plant-rooted-life:surfaces}{根毛}为土壤提供了广阔的表面，
以及茎叶系统，其叶片为光和
空气提供了广阔的表面。
\end{solution}

\begin{solution}{exo:g12:plant-rooted-life:2}
叶表皮上由两个保卫\omterm{def:g10:cells-common-unit:cell}{细胞}围成的孔隙。二氧化碳
进入，水蒸气（及氧气）离开。光下张开，
夜间和干旱时关闭。
\end{solution}

\begin{solution}{exo:g12:plant-rooted-life:3}
原汁：水和无机离子，在\omterm{prop:g12:plant-rooted-life:saps}{木质部}中由根运往叶，
由蒸腾作用拉动。筛管汁液：水和\omterm{def:g10:chemistry-of-life:families}{糖类}，在
\omterm{prop:g12:plant-rooted-life:saps}{韧皮部}中，由叶运往利用或储存糖的器官，
双向运输，由装载了糖的\omterm{def:g10:cells-common-unit:cell}{细胞}所产生的压强推动。
\end{solution}

\begin{solution}{exo:g12:plant-rooted-life:4}
结构防御：角质层和树皮、刺（还有刺毛、二氧化硅、
\omterm{def:g10:cells-common-unit:organelle}{细胞壁}）。化学防御：单宁、生物碱（还有毒素、树脂、
乳汁）。
\end{solution}

\begin{solution}{exo:g12:plant-rooted-life:5}
对刺激（光、重力、水）的定向生长反应。
生长素，由茎的生长尖端合成。
\end{solution}

\begin{solution}{exo:g12:plant-rooted-life:6}
$\qty{40}{cm^2} = \qty{4000}{mm^2}$: $4000 \times 200 = \num{8e5}$
每片叶的\omterm{def:g12:plant-rooted-life:stomata}{气孔}数；全树约 $\num{1.6e11}$ 个。
\end{solution}

\begin{solution}{exo:g12:plant-rooted-life:7}
约在正午至午后早些时候。午后炎热时，保卫\omterm{def:g10:cells-common-unit:cell}{细胞}
部分关闭\omterm{def:g12:plant-rooted-life:stomata}{气孔}以减少失水，并且
蒸腾作用随\omterm{def:g12:plant-rooted-life:stomata}{气孔}的开闭而变化，而非随光照。
\end{solution}

\begin{solution}{exo:g12:plant-rooted-life:8}
叶片保持绿色，水分供应充足（\omterm{prop:g12:plant-rooted-life:saps}{木质部}位于
木材内部，完好无损）。根被切断了叶片提供的糖，
其储备在数月内耗尽而死亡，整棵树随之死去。紧贴环剥处上方，
树皮因无法下运的糖分而膨大。
\end{solution}

\begin{solution}{exo:g12:plant-rooted-life:9}
完整时：存在反应。尖端遮盖：尖端是感受光的部位，
基部遮盖：基部弯曲无需受光，它响应的是
信号，而非光。去除尖端：尖端不可或缺，
信号由其产生。三组实验合起来，把感受（尖端）与
反应（基部）分开，并表明两者之间有某种物质传递。
\end{solution}

\begin{solution}{exo:g12:plant-rooted-life:10}
能量来自太阳：叶表面的水分蒸发，由太阳热量
和干燥空气驱动，拉动\omterm{prop:g12:plant-rooted-life:saps}{木质部}中连续的水柱向上。
植物提供的是导管和张开的\omterm{def:g12:plant-rooted-life:stomata}{气孔}，而非
能量。
\end{solution}

\begin{solution}{exo:g12:plant-rooted-life:11}
风媒花：小、不起眼、无香味、无花蜜，产生大量轻
花粉，柱头呈羽毛状——风无定向且无需报酬，因此植物
以数量取胜。虫媒花：大、色彩鲜艳、有香味、
有花蜜，花粉黏性强、数量适中——传粉者必须被
吸引并给予回报，并能准确传递花粉。
\end{solution}

\begin{solution}{exo:g12:plant-rooted-life:12}
它获得水分，使\omterm{def:g10:cells-common-unit:cell}{细胞}保持膨压而存活；却失去了
二氧化碳的进入，因而无法进行\omterm{def:g10:cell-metabolism:autotrophy}{光合作用}，也失去了
蒸腾作用带来的降温。\omterm{def:g12:plant-rooted-life:stomata}{气孔}关闭意味着无糖合成，植物依靠
自身储备维生；长期干旱会先使其饥饿，而非先使其失水。
\end{solution}

\begin{solution}{exo:g12:plant-rooted-life:13}
在叶中起防御作用：一种针对植食性动物的毒素，嵌入
组织之中。在花蜜中起招引作用：低剂量使
传粉者记住并返回，与糖一同作为回报。同一种
物质满足两种需要，剂量随部位而异。
\end{solution}

\begin{solution}{exo:g12:plant-rooted-life:14}
花管稍长的花只被能够
探到底部的蛾类传粉，因此其花粉只传给同类花；
口器稍长的蛾类能获得他类无法获取的花蜜；每一项
微小优势都会被选择，花管与口器在
多代之间共同延长。若蛾类消失，花便
不能结实；若花消失，蛾类便失去唯一的食物。
\end{solution}

\begin{solution}{exo:g12:plant-rooted-life:15}
植物用茎尖感受光，用根感受重力；
它分别以向光生长和向重力生长作出反应；它
以树皮和刺自卫，并在受攻击后数小时内产生毒素和消化
抑制剂。它还招引昆虫搬运花粉，招引鸟类搬运
种子，招引黄蜂消灭其毛毛虫，
并以花蜜、果实和气味作为回报。它无需移
动即可行动。
\end{solution}

\begin{solution}{pb:g12:plant-rooted-life:1}
\textbf{1.} $\num{250000} \times \qty{25}{cm^2} = \num{6.25e6}\,
\unit{cm^2} = \qty{625}{m^2}$.

\textbf{2.} $\qty{625}{m^2} = \num{6.25e8}\,\unit{mm^2}$；$\times 250
\approx \num{1.6e11}$ 个\omterm{def:g12:plant-rooted-life:stomata}{气孔}。

\textbf{3.} 每平方米地面上约有 $625/150 \approx 4$ 平方米的叶片，
即叶面积约为地面面积的 4 倍。

\textbf{4.} \omterm{prop:g12:plant-rooted-life:surfaces}{根毛}面积约为 \qty{1000}{m^2}，与 \qty{625}{m^2}
的叶面积相当：两者都是与相同物质流相匹配的交换表面，
即叶片散失的水必须由根吸收。

\textbf{5.} 光在组织表层不到一毫米的范围内即被吸收，而二氧化碳
只能扩散很短的距离：若器官较厚，
则大部分\omterm{def:g10:cells-common-unit:cell}{细胞}将处于黑暗之中并陷入饥饿状态。
叶片薄且层叠排列，使每个\omterm{def:g10:cells-common-unit:cell}{细胞}都靠近光与空气。

\textbf{6.} $625 \times 0.6 = \qty{375}{L}$。

\textbf{7.} 二氧化碳 $625 \times 8 = \qty{5}{kg}$，可得
约 $5/1.47 \approx \qty{3.4}{kg}$ 的糖（或简言之：
由 \qty{5}{kg} 的 $\mathrm{CO_2}$ 合成的糖约为 \qty{3.4}{kg}）；
耗水量为 $3.4/1.6 \approx \qty{2.1}{kg}$，约占
蒸腾水量 \qty{375}{L} 的 0.6\%。

\textbf{8.} 二氧化碳经由让水分逸出的同一批开放\omterm{def:g12:plant-rooted-life:stomata}{气孔}进入叶片；
关闭\omterm{def:g12:plant-rooted-life:stomata}{气孔}则\omterm{def:g10:cell-metabolism:autotrophy}{光合作用}停止。水分是
合成糖所付出的代价。

\textbf{9.} 动力来自太阳，它使叶片上的水分蒸发。将
\qty{375}{kg} 的水提升 \qty{20}{m}，约需 \qty{75}{kJ} 的功——本身并不多，
但若要用生物泵完成此过程，则需要组织、ATP 以及
树木无须额外构建的维持机制。

\textbf{10.} 为抵御干热，它已部分关闭\omterm{def:g12:plant-rooted-life:stomata}{气孔}；
二氧化碳的进入量与糖的产量随之按相近比例下降，
——整个下午约损失 \qty{2}{kg} 糖。

\textbf{11.} 约制成 \qty{3.4}{kg} 糖；其中 \qty{2}{kg}（60\%）输送至根、
树干和嫩枝，\qty{0.5}{kg}（15\%）输送至橡实。

\textbf{12.} \omterm{prop:g12:plant-rooted-life:saps}{韧皮部}把糖从叶片（源）运至根、树干、
嫩枝和橡实（库）。

\textbf{13.} 以淀粉形式储存于树干和根中；次年春季，
它再转化为糖，用以建造新叶，
使新叶在能自行制造养分之前得以生长。

\textbf{14.} 环割线以上的嫩枝和橡实仍能获得糖；
根和树干下部则一无所得：其原本的 \qty{2}{kg} 份额
降至储存量所能提供的水平，它们将在这一季中逐渐饿死。

\textbf{15.} 其余叶片仍持续制糖，树干和根中储存的淀粉
可弥补亏缺，树木还能从众多生长点
长出新叶；若夏季的糖量足够，橡实这一库
仍可被充实。

\textbf{16.} 花小而呈绿色，无花瓣、无香气、无花蜜，
成串下垂为柔荑花序，释放大量轻质花粉；
另有大而羽毛状的柱头用于接收花粉。

\textbf{17.} 橡实若落在母树下方，即便能萌发，也只能
生长于其荫蔽之中并与之竞争；松鸦可将其搬运数百米，
并埋入适宜深度。丢失其中大部分，
是为将少数种子安置得当所付出的代价。

\textbf{18.} $3000 \times 0.05 = 150$ 颗被遗忘的橡实；$150/50 = 3$
株幼树。

\textbf{19.} 在数十年的丰收年份里，每年约产生 3 株幼树，
树木一生中或可产生约百株幼树，其中须有一株存活至成体。
数量之所以如此之大，是因为几乎每一颗种子和幼树
都会遭到损失。

\textbf{20.} \qty{625}{m^2} 的叶片，约 $\num{1.6e11}$ 个\omterm{def:g12:plant-rooted-life:stomata}{气孔}；
约 \qty{375}{L} 的水在一天内被提升，以换取 \qty{3.4}{kg} 糖；
借风传粉，借松鸦传播橡实。
\end{solution}
